\documentclass[11pt,a4paper]{article}
\usepackage[margin=20mm,headheight=14pt]{geometry}
\usepackage[T1]{fontenc}
\usepackage[utf8]{inputenc}
\usepackage{lmodern,microtype,amsmath,amssymb,array,booktabs,xcolor,fancyhdr,enumitem}
\usepackage[most]{tcolorbox}
\definecolor{ink}{HTML}{17334C}
\definecolor{pale}{HTML}{EEF4F7}
\definecolor{rust}{HTML}{A64B2A}
\definecolor{warn}{HTML}{8A3324}
\definecolor{palewarn}{HTML}{FBEEE9}
\setlength{\parindent}{0pt}
\setlength{\parskip}{6pt}
\pagestyle{fancy}\fancyhf{}
\fancyhead[L]{\small\color{ink} O RACIOCÍNIO E A MATEMÁTICA}
\fancyhead[R]{\small 28.09.2026}
\fancyfoot[C]{\small\thepage}
\renewcommand{\headrulewidth}{0.3pt}
\newtcolorbox{caixa}[1]{colback=pale,colframe=ink,boxrule=0.4pt,arc=1mm,
 title={#1},fonttitle=\bfseries,coltitle=white,left=3mm,right=3mm,top=2mm,bottom=2mm}
\newtcolorbox{alerta}[1]{colback=palewarn,colframe=warn,boxrule=0.4pt,arc=1mm,
 title={#1},fonttitle=\bfseries,coltitle=white,left=3mm,right=3mm,top=2mm,bottom=2mm}
\newcommand{\Mst}{M_*}

\begin{document}

{\Large\bfseries\color{ink} A matemática e o raciocínio que eu segui}\par
{\large Do observável de simetria de 1-forma até onde o argumento trava}\par
\vspace{1mm}
{\small Mauro de Oliveira Cardoso (Viole) --- 28 de setembro de 2026.
Acompanha o artigo canônico v0.3 e o resumo de seis páginas.}

\vspace{2mm}\hrule\vspace{2mm}

\section*{Passo 1 --- De onde eu parto}

Numa teoria de Maxwell, a simetria elétrica de 1-forma mede carga através de
uma superfície que envolve uma linha de prova. Com matéria carregada, o
\emph{screening} faz essa carga depender do raio. Córdova, Ohmori e Rudelius
definiram o perfil que mede essa quebra; Basile e Golmohammadi calcularam-no a
um laço:
\[
  \delta(r) \;=\; -\frac{r}{q_\infty}\,\frac{dq(r)}{dr},
  \qquad q_\infty=\lim_{r\to\infty}q(r)\neq 0 .
\]
\textbf{Eu não calculo esse perfil.} Isso já estava feito. Eu pergunto outra
coisa.

\begin{caixa}{A MINHA PERGUNTA}
O \emph{formato} de $\delta(r)$ diz alguma coisa sobre o espectro da teoria?
\end{caixa}

\section*{Passo 2 --- A ideia central, numa linha}

Defino o perfil reduzido $\Phi(r)=\delta(r)/r^2$. A afirmação é:

\begin{caixa}{TEOREMA A (condicional)}
Sob as hipóteses H1--H5, e em particular sob \textbf{H3},
\[
  \Phi(r)\;=\;\int_0^\infty e^{-rx}\,d\nu(x),\qquad \nu\ \ge\ 0,
\]
\emph{exatamente} --- não só em ordem líder. A borda do suporte de $\nu$,
$\Mst=\inf\operatorname{supp}\nu$, é o menor limiar que \textbf{acopla a este
observável}.
\end{caixa}

O mecanismo é uma identidade de núcleo: derivando
$(1+r\sqrt{s})e^{-r\sqrt{s}}$ em $r$ sai exatamente $-s\,r\,e^{-r\sqrt{s}}$.
O fator de Gauss cancela e sobra uma transformada de Laplace pura.

\section*{Passo 3 --- Por que isso é útil}

Porque transformada de Laplace de medida positiva é objeto clássico, com
ferramentas prontas:

\begin{enumerate}[leftmargin=7mm,itemsep=3pt]
\item \textbf{Inclinação logarítmica.} $\Gamma(r)=-\partial_r\log\Phi(r)
  =\langle x\rangle_r \ \ge\ \Mst$, e $\Gamma$ decresce em $r$. É a ``massa
  efetiva'' da rede, numa variável radial.
\item \textbf{Hierarquia de Hankel / GEVP.} Com os momentos
  $a_n(r)=(-1)^n\Phi^{(n)}(r)$, monto $H_0=(a_{i+j})$, $H_1=(a_{i+j+1})$ e
  tomo $B_K(r)=\lambda_{\min}(H_1,H_0)\ \ge\ \Mst$, com
  $B_{K+1}\le B_K$ e $B_0=\Gamma$.
\item \textbf{Lei de borda.} Se a densidade vai como $C t^{p-1}$ perto da
  borda ($t=x-\Mst$), então
  $\Gamma(r)=\Mst+p/r+\beta p/r^2+O(r^{-3})$. Para um Dirac, $p=3/2$ e
  $\Gamma=2m+\tfrac{3}{2r}-\tfrac{5}{16mr^2}+\cdots$
\end{enumerate}

\textbf{Todas dão cotas superiores que descem até $\Mst$.} A matemática aqui é
clássica; eu não reivindico novidade nela.

\section*{Passo 4 --- Onde eu fui honesto sobre os limites}

\textbf{(a) O limiar não é a massa do elétron.} $x=\sqrt{s}$ é massa invariante
do estado intermediário. A um laço a borda é $2m$ --- um \emph{par}, não uma
partícula.

\textbf{(b) O limiar pode nem ser de matéria carregada.} Em QED completa há
cortes multi-fóton desde $s=0$. O gap some e $\Gamma$ devolve o ínfimo do
suporte \emph{efetivamente acoplado}, que pode ser zero. Por isso eu escrevo
``limiar acoplado'', não ``limiar carregado''.

\textbf{(c) Não existe cota inferior uniforme (Teorema H).} Este é o resultado
que eu mais gosto porque não depende de H3. Tome
\[
  f_0=e^{-Mt},\qquad f_\epsilon=(1-\epsilon)e^{-Mt}+\epsilon e^{-\mu t},
  \qquad \mu<M .
\]
As duas são Laplace de medidas positivas, com bordas $M$ e $\mu$, mas
$\sup_{[0,T]}|f_\epsilon-f_0|\le\epsilon$. Com dados de precisão finita numa
janela finita, \textbf{nenhuma cota inferior estritamente positiva} para $\Mst$
sobrevive. A trivial $\Mst\ge0$ permanece.

\begin{caixa}{O EXEMPLO CONCRETO --- ``átomo minúsculo''}
$d\nu=10^{-12}\delta_2+(x-3)^{1/2}e^{-(x-3)}\mathbf{1}_{x\ge3}dx$.
Até $r\approx 22{,}76$ o perfil é indistinguível de um espectro que começa em 3.
A inclinação só cai para 2 depois disso.
\textbf{As cotas continuam corretas --- só deixam de ser informativas.}
\end{caixa}

\textbf{(d) Positividade sozinha não dá a conjectura da gravidade fraca.}
Há uma obstrução de escala: a positividade é invariante por reescalonamento, e
a conjectura não é. A implicação que eu enuncio é condicional e exige insumo
gravitacional externo, declarado antes.

\section*{Passo 5 --- A hipótese H3, que é onde tudo trava}

H3 diz que o kernel de resposta estática é uma função de Stieltjes:
\[
  \mathcal{G}(Q^2)\;=\;\frac{g_R^2}{Q^2\bigl[1-g_R^2\overline{\Pi}(Q^2)\bigr]},
  \qquad
  \overline{\Pi}(Q^2)=Q^2\!\int\!\frac{\rho_J(s)\,ds}{s(s+Q^2)}\ \ge\ 0 .
\]

\textbf{Em ordem líder, H3 sai de graça.} A positividade de Lehmann da corrente
de matéria dá $\rho_J\ge0$, e daí
$d\sigma_{\rm LO}=g_R^4\,\rho_J(s)\,ds/s\ \ge 0$. É o cálculo de Uehling; o
coeficiente $2\alpha/3\pi$ é reproduzido.

\begin{alerta}{A ORDEM SEGUINTE NÃO TEM A FORMA CERTA --- e isso importa}
O coeficiente redutível de $O(g_R^6)$ é
$g_R^6\,Q^2\bigl[\int\rho_J/(s(s+Q^2))\bigr]^2$. Essa função vale $0$ em
$Q^2=0$, cresce, e decai como $1/Q^2$: tem \textbf{máximo interior}.
Mas toda $\int d\sigma/(Q^2+s)$ com $\sigma\ge0$ e suporte longe de zero é
\textbf{positiva em $Q^2=0$ e estritamente decrescente}.
Logo o coeficiente de ordem fixa \textbf{não} é Stieltjes-positivo.
\end{alerta}

A saída é que a positividade só aparece no objeto \textbf{ressomado}. No modelo
solúvel $\rho_J=Z\delta(s-s_0)$, com $c=g_R^2Z/s_0\in(0,1)$:
\[
  \mathcal{G}=\frac{g_R^2}{Q^2}
   +\frac{g_R^2c/(1-c)}{Q^2+s_0/(1-c)} .
\]
O peso é positivo. O que acontece é que o \textbf{polo se desloca} com o
acoplamento. Expandir em $g_R$ a $s$ fixo produz coeficientes de sinal
alternado --- daí a aparente violação ordem a ordem.

\section*{Passo 6 --- O que eu achei nesta rodada: o critério $Z_3$}

Por frações parciais, $\dfrac{Q^2}{s(s+Q^2)}=\dfrac1s-\dfrac1{s+Q^2}$, logo
\[
  \boxed{\;1-g_R^2\overline{\Pi}(Q^2)
   \;=\;\underbrace{\bigl[1-g_R^2\overline{\Pi}(\infty)\bigr]}_{\textstyle =:Z_3}
   \;+\;g_R^2\!\int\!\frac{\rho_J(s)\,ds}{s+Q^2}\;}
\]
O lado direito é \emph{exatamente} a forma de Stieltjes. Usando a dualidade
clássica (Schilling--Song--Vondra\v{c}ek, Teor.~7.3: $f\not\equiv0$ é Stieltjes
$\iff$ $1/f$ é Bernstein completa $\iff$ $z f(z)$ é Bernstein completa),
segue o critério.

\begin{alerta}{CORREÇÃO DE FRONTEIRA (28/09) --- enunciar assim, não mais forte}
\textbf{(a)} $Z_3\ge0$ caracteriza a classe \textbf{geral} de Stieltjes, que
admite uma \emph{constante aditiva}.\\
\textbf{(b)} A H3 do artigo é escrita \textbf{sem} essa constante, logo exige
$\mathcal{G}\to0$. Aí vale só a \textbf{suficiência}: $Z_3>0$ dá
$0<\mathcal{G}\le g_R^2/(Q^2Z_3)\to0$.\\
\textbf{(c)} Na fronteira $Z_3=0$ não falha sempre: $\mathcal{G}\to1/\rho((0,\infty))$,
então falha \emph{se e somente se} $\rho$ tem massa total finita.
Exemplo de falha: $\rho=Z\delta(s-s_0)$ dá $\mathcal{G}=g_R^2/Q^2+g_R^2/s_0$.
\end{alerta}

\textbf{O que o critério diz sobre QED de verdade} (um laço, $\alpha=1/137$,
unidades de $m_e$):

\begin{center}\small
\begin{tabular}{lrrl}
\toprule
$\Lambda/m$ & $\log(\Lambda^2/4m^2)$ & $Z_3$ & H3? \\
\midrule
$10^{6}$    & 26{,}2   & 0{,}980  & vale \\
$10^{19}$   & 86{,}1   & 0{,}934  & vale (escala de Planck) \\
$10^{100}$  & 459      & 0{,}645  & vale \\
$e^{645}$   & 1288{,}6 & 0{,}0022 & vale, no limite \\
$e^{700}$   & 1398{,}6 & $-0{,}083$ & \textbf{falha} \\
\bottomrule
\end{tabular}
\end{center}

O polo de Landau está em $\log(\Lambda^2/4m^2)=3\pi/\alpha\approx1291$.
\textbf{Para todo cutoff fisicamente significativo, H3 vale com folga enorme.}
Falha só no contínuo estrito --- e eu não escondo isso.

\section*{Passo 6b --- O mecanismo, a fronteira, e o que os momentos n\~ao veem}

\textbf{Tr\^es regimes, n\~ao dois.} Como $W(0)=1$ e $W(\infty)=Z_3$, com $W$
estritamente decrescente:

\begin{center}\small
\begin{tabular}{llll}
\toprule
Regime & Polo spacelike & \'Atomo acima do corte & H3 (sem constante) \\
\midrule
$Z_3>0$ & n\~ao existe & \textbf{sim}, um, peso $>0$ & \textbf{vale} \\
$Z_3=0$ & n\~ao existe & --- & \textbf{falha} se a massa total for finita \\
$Z_3<0$ & existe, res\'iduo $<0$ & n\~ao existe & \textbf{falha} \\
\bottomrule
\end{tabular}
\end{center}

O polo em $Q^2>0$ \'e um estado com $s<0$ --- taquiônico. Uma fun\c{c}\~ao de
Stieltjes n\~ao pode t\^e-lo. \textbf{Mas ausência de polo n\~ao basta:} na
fronteira $Z_3=0$ n\~ao h\'a polo e H3 ainda falha, por um termo constante.
A aus\^encia de polo \'e \textbf{necess\'aria, n\~ao suficiente}.

\begin{caixa}{COMO EU SEI QUE N\~AO FALTA TERMO ESPECTRAL \emph{NESTE MODELO}}
Depois de achar o \'atomo esquecido, a pergunta \'obvia \'e: falta outro?

\textbf{Escopo da resposta:} ela vale para o \emph{modelo RPA} com as hip\'oteses
declaradas (densidade $\ge0$ e n\~ao nula, suporte positivo, momento inverso
finito, acoplamento positivo, representa\c{c}\~ao de Dyson subtra\'ida), e \'e sobre
\textbf{completude da representa\c{c}\~ao}, n\~ao sobre precis\~ao da implementa\c{c}\~ao.

\textbf{N\~ao falta termo.} $W$ n\~ao tem zero fora do eixo real, porque
$\operatorname{Im}W(x+iy)=-g_R^2\,y\int\rho_J/|s+z|^2$ tem o sinal de $-y$ e
n\~ao se anula. Ent\~ao todos os zeros s\~ao reais, e o eixo real tem quatro
regi\~oes, todas conferidas: spacelike, abaixo do limiar, o corte, e acima do
corte. Nenhuma outra contribui\c{c}\~ao pode existir.

\textbf{E h\'a um controle global.} Duas \emph{regras de soma} obrigam a somar
\emph{todo} o peso espectral:
\[g_R^2+\int d\sigma+\sum_a w_a=\frac{g_R^2}{Z_3},\qquad
\int s\,d\sigma+\sum_a w_a s_a=\frac{g_R^4\mu}{Z_3^2}.\]
Conferem a $6\times10^{-10}$ e $5\times10^{-9}$. Sem o \'atomo, a primeira
piora para $6\times10^{-5}$. \'E o teste que pega peso esquecido --- concord\^ancia
em alguns pontos, n\~ao.

\textbf{Sobre precis\~ao, separadamente.} Comparei o caminho r\'apido com
aritm\'etica de multiprecis\~ao e com um caso de forma fechada. O erro real \'e
$\sim2\times10^{-6}$ na posi\c{c}\~ao e no peso do \'atomo --- bem maior que a
toler\^ancia pedida \`a quadratura, porque o integrando tem ponto de
ramifica\c{c}\~ao no extremo. \textbf{Completude e precis\~ao s\~ao coisas
diferentes}, e eu s\'o sei afirmar a primeira por demonstra\c{c}\~ao.
\end{caixa}

\begin{alerta}{ERRO MEU, CORRIGIDO EM 29/09 --- vale contar se ele perguntar}
Eu havia afirmado que a medida espectral fica dentro do suporte da densidade de
entrada. \textbf{\'E falso.} A ressoma cria um \textbf{\'atomo discreto acima do
corte}: aqui em $s_a=1{,}0000053\times10^6$ com peso $6{,}27\times10^{-4}>0$.

Pior: a concord\^ancia de $10^{-8}$ que eu atribu\'ia a erro de quadratura
\textbf{era o \'atomo omitido}. Inclu\'i-lo melhora a reconstru\c{c}\~ao em
\textbf{quatro ordens de grandeza}. A positividade N\~AO foi afetada --- o peso
\'e positivo ---, mas a descri\c{c}\~ao do suporte estava errada.

O \'atomo \'e artefato do corte r\'igido: com regulador suave n\~ao sobra
regi\~ao real onde ele possa aparecer.
\end{alerta}

\begin{alerta}{O QUE MAIS ME SURPREENDEU --- e eu levo isto ao professor}
Eu esperava que a viola\c{c}\~ao aparecesse como densidade espectral negativa.
\textbf{N\~ao aparece.} A densidade do cont\'inuo \'e
$d\sigma/ds=g_R^4\rho_J/(s|W|^2)$, manifestamente $\ge0$ para
\emph{qualquer} acoplamento --- \'e $\rho_J$ dividida por um m\'odulo ao quadrado.

Consequ\^encia: um teste de positividade baseado nos \textbf{momentos do
cont\'inuo} aprova o caso mesmo quando H3 \'e falsa. Rodei o filtro do pr\'oprio
projeto e ele devolve \textsc{checked\_compatible} nos dois lados da fronteira.
N\~ao \'e defeito de implementa\c{c}\~ao: nenhum momento do cont\'inuo enxerga um
polo discreto. A viola\c{c}\~ao inteira mora no polo.
\end{alerta}

\textbf{Evid\^encia computacional.} 211 testes passando, tamb\'em em
\emph{checkout} limpo do reposit\'orio, com os geradores reproduzindo os
artefatos byte a byte. As quatro condi\c{c}\~oes que identificam a
representa\c{c}\~ao do modelo com a forma exata de H3 do artigo --- res\'iduo de
Coulomb $g_R^2>0$, suporte fora da origem, aus\^encia de constante e momento
inverso finito --- foram verificadas; as tr\^es primeiras demonstradas, a quarta
demonstrada para a densidade concreta e verificada numericamente.
Os erros de quadratura s\~ao estimativas do QUADPACK: \textbf{evid\^encia
num\'erica}, n\~ao certifica\c{c}\~ao por intervalo.

\section*{Passo 7 --- Onde o argumento realmente trava}

O critério acima é sobre a \textbf{cadeia de Dyson} com $\rho_J$ vindo da
corrente de matéria \emph{não gaugeada}, onde a positividade de Lehmann é
limpa. Isso \textbf{não} é a resposta estática não perturbativa e invariante de
gauge. A rota séria de teoria de campos é outra:

\begin{enumerate}[leftmargin=7mm,itemsep=3pt]
\item Com axiomas de Wightman e identidade de Bianchi, \textbf{sem} supor
  $F=dA$, o correlator $\langle F_{\mu\nu}F_{\rho\sigma}\rangle$ tem
  representação com medida positiva. Bianchi elimina a estrutura dual para
  $p^2>0$; em $p^2=0$ \textbf{não} elimina, e sobra um termo que só a
  localidade, via CPT, remove. \emph{Esse lema eu ainda não li na fonte.}
\item Transportar essa positividade até o laço de Wilson, e daí ao kernel
  estático. \textbf{Este passo não está feito.} Envolve projeção tensorial,
  termos de contato, renormalização de perímetro e o limite $T\to\infty$.
\end{enumerate}

Um pedaço eu verifiquei: no limite estático,
$\langle E_iE_j\rangle=D(Q^2)\,p_ip_j$ é \textbf{puramente longitudinal}, e o
escalar que o potencial vê é $Q^2D(Q^2)=g^2_{\rm eff}$. Isso dá uma
reformulação útil: $\mathcal{G}$ está na classe geral de Stieltjes
$\iff$ $g^2_{\rm eff}$ é Bernstein completa --- e $g^2_{\rm eff}$ é
exatamente o que o correlator entrega, sem inversão nenhuma.

\begin{caixa}{A PERGUNTA QUE EU LEVO AO PROFESSOR}
Esse transporte --- do correlator $\langle FF\rangle$ para o kernel do laço de
Wilson --- esconde alguma hipótese além da fase de Coulomb?
\end{caixa}

\section*{Passo 8 --- O que continua em aberto}

\begin{itemize}[leftmargin=6mm,itemsep=2pt]
\item H3 além da ordem líder para o kernel interagente: \textbf{não provada}.
\item O transporte $\langle FF\rangle\to$ Wilson $\to$ kernel: \textbf{não feito}.
\item Fase de Coulomb tem de ser hipótese explícita --- Proca satisfaz os
  axiomas e \emph{não} tem polo de Coulomb ($q_\infty=0$, $\delta$ nem definida).
\item Isolar canais carregados dos cortes sem massa preservando positividade.
\item Critério $Z_3$: nota de trabalho, não promovida ao artigo; atribuição em
  aberto (Brown--Weisberger 1979 não lido; Källén 1952 é clássico).
\item Nenhum CI remoto executado; nenhuma revisão humana do PDF; nada submetido.
\end{itemize}

\vspace{2mm}\hrule\vspace{2mm}
{\small\color{rust}\textbf{Resumo em uma frase.} Sob uma hipótese de
positividade que vale em ordem líder e que eu sei caracterizar na cadeia de
Dyson, o formato do perfil determina cotas superiores para o limiar acoplado ---
e eu sei exatamente onde o argumento ainda não fecha.}

\end{document}
